\begin{figure}[htbp]
  \centering
  \textbf{\small (a) Ficha da referência --- proposta, não implementada}\\[3pt]
  \begin{tikzpicture}[x=1cm, y=1cm, font=\scriptsize,
      cmp/.style={draw=corCinzaC, line width=.5pt, fill=white, align=left,
                  inner sep=3pt, font=\scriptsize, text=corTinta,
                  text width=3.9cm, minimum height=13mm, anchor=north west}]
    \providecommand{\rot}{}\renewcommand{\rot}[1]{{\tiny\color{corCinza}#1}}
    \begin{scope}[on background layer]
      \fill[corFundo, rounded corners=2pt] (0,-3.95) rectangle (13.6,1.30);
      \draw[corCinza, line width=.6pt, dash pattern=on 3pt off 2pt,
            rounded corners=2pt] (0,-3.95) rectangle (13.6,1.30);
    \end{scope}
    \node[anchor=north west, font=\scriptsize\bfseries, text=corTinta] at (0.30,1.15)
      {Destilação atmosférica --- fuel gás};
    \node[anchor=north east, font=\tiny, text=corCinza] at (13.30,1.12)
      {referência v4, de 2026-03-02 \quad$\vert$\quad versão anterior consultável};
    \draw[corCinzaC, line width=.5pt] (0.30,0.62) -- (13.30,0.62);

    \node[cmp] at (0.30,0.42)
      {\rot{DOMÍNIO DE VALIDADE}\\ carga entre 18,6 e 27,4~kt/d\\
       \rot{fora do domínio não há juízo}};
    \node[cmp] at (4.80,0.42)
      {\rot{DESVIO DO DIA}\\ $+$41 t GNE \quad($+$12,8\,\%)\\
       \rot{observado $-$ esperado}};
    \node[cmp] at (9.30,0.42)
      {\rot{INCERTEZA}\\ banda de 95\,\%: $\pm$94 t GNE\\
       \rot{cobertura verificada: 0,87}};

    \node[cmp, draw=corFG, line width=.7pt, fill=corFG!7] at (0.30,-1.20)
      {\rot{ESTADO}\\ \textbf{em observação}\\
       \rot{desde 2026-07-14, 4.\textordmasculine{} dia}};
    \node[cmp] at (4.80,-1.20)
      {\rot{AÇÃO EM CURSO}\\ verificar medição do canal\\
       \rot{responsável: engenharia de processo}};
    \node[cmp] at (9.30,-1.20)
      {\rot{QUALIDADE DO DIA}\\ 3 origens concordantes\\
       \rot{sem imputação}};

    \draw[corCinzaC, line width=.5pt] (0.30,-2.85) -- (13.30,-2.85);
    \node[anchor=north west, font=\tiny, text=corCinza, text width=13.0cm]
      at (0.30,-2.98)
      {Seis campos na ficha principal. Os restantes campos do registo ---
       parâmetros, histórico de versões, diagnóstico completo --- ficam no
       registo detalhado (Apêndice~\ref{app:manual}) e não competem com estes
       pela atenção de quem decide.};
  \end{tikzpicture}

  \vspace{4mm}
  \textbf{\small (b) Estados da referência e eventos que os mudam}\\[3pt]
  \begin{tikzpicture}[x=1cm, y=1cm, font=\scriptsize,
      st/.style={ntcaixa, text width=2.3cm, minimum height=10mm,
                 draw=corFG, line width=.7pt, fill=corFG!7},
      ev/.style={font=\tiny, text=corCinza, align=center, text width=1.45cm,
                 inner sep=1.5pt, fill=white}]
    \node[st] (op) at (1.45,0) {\textbf{em vigor}};
    \node[st] (ob) at (5.35,0) {\textbf{em observação}};
    \node[st] (su) at (9.25,0) {\textbf{suspensa}};
    \node[st] (ve) at (13.15,0) {\textbf{em verificação}};

    \draw[ntsetaf] (op) -- (ob);
    \node[ev, anchor=south] at (3.40,0.12) {sinal do monitor};
    \draw[ntsetaf] (ob) -- (su);
    \node[ev, anchor=south] at (7.30,0.12) {causa confirmada ou domínio violado};
    \draw[ntsetaf] (su) -- (ve);
    \node[ev, anchor=south] at (11.20,0.12) {reestimação datada};

    \draw[ntsetaf] (ve.north) to[out=105, in=75, looseness=0.55] (op.north);
    \node[ev, anchor=south, text width=5.0cm] at (7.30,1.55)
      {cobertura verificada fora da amostra};
    \draw[ntsetaf] (ob.south) to[out=-110, in=-70, looseness=0.9] (op.south);
    \node[ev, anchor=north, text width=4.2cm] at (3.40,-1.34)
      {sinal não confirmado em 10 dias};

    \node[ntmini, anchor=north west, text width=13.5cm] at (0.20,-2.10)
      {Os estados persistem e ficam datados; os eventos que os mudam são
       instantâneos e ficam registados com quem os desencadeou. Uma
       referência suspensa não produz juízo: a unidade continua a ser
       medida, mas não é julgada até a nova referência passar a verificação.};
  \end{tikzpicture}
  \caption[Ficha da referência e diagrama de estados]%
  {Proposta de ficha da referência e do ciclo de vida que a governa. Os
   valores mostrados são ilustrativos. O painel~(a) mostra o que quem
   decide precisa de ver de uma vez: onde a referência é válida, qual o
   desvio, com que incerteza, em que estado está, o que se está a fazer e
   que qualidade tem o dado do dia. O painel~(b) separa aquilo que o método
   em produção não separa: o alerta é um evento, o estado da referência é
   uma condição persistente, e a passagem de um ao outro exige uma decisão
   com dono e data. Nada disto foi implementado neste trabalho.}
  \label{fig:ficha-estados}
\end{figure}
